\subsection{Calculus Vol2 - 5.5 Alternating Series}

\url{https://openstax.org/books/calculus-volume-2/pages/5-5-alternating-series}

\begin{enumerate}
  \item What is an alternating series? Write its general form.
  \item How can an alternating series be written using the factor $(-1)^n$ or $(-1)^{n+1}$?
  \item In the representation $\sum_{n=1}^{\infty} (-1)^{n+1} b_n$, what conditions are imposed on $b_n$?
  \item State the Alternating Series Test (Leibniz Test).
  \item What two conditions must the sequence $\{b_n\}$ satisfy for the Alternating Series Test to apply?
  \item Why is the condition $\lim_{n \to \infty} b_n = 0$ necessary for convergence?
  \item Does the Alternating Series Test guarantee absolute convergence or only convergence? Explain.
  \item Give an example of a series that converges by the Alternating Series Test.
  \item What is meant by absolute convergence?
  \item What is meant by conditional convergence?
  \item What is the relationship between absolute convergence and convergence of a series?
  \item Give an example of a series that converges conditionally but not absolutely.
  \item How can you test whether an alternating series converges absolutely?
  \item If $\sum |a_n|$ diverges but $\sum a_n$ converges, what type of convergence does the series have?
  \item What is the alternating harmonic series? Does it converge absolutely or conditionally?
  \item Define the nth partial sum $S_N$ of a series.
  \item What is the remainder $R_N$ of a series?
  \item For an alternating series satisfying the Alternating Series Test, what inequality bounds the remainder $R_N$?
  \item Write the inequality relating $|R_N|$ and $b_{N+1}$.
  \item What does the remainder estimate tell us about the error when approximating a sum by $S_N$?
  \item How can you use the remainder estimate to determine how many terms are needed for a desired accuracy?
  \item If you approximate the sum of an alternating series by $S_N$, what is the maximum possible error?
  \item Why do partial sums of an alternating series tend to oscillate around the true sum?
  \item Describe the behavior of even and odd partial sums in an alternating series that satisfies the test.
  \item How can you estimate the sum of an alternating series to within a given tolerance $\varepsilon$?
  \item For the series $\sum_{n=1}^{\infty} \frac{(-1)^{n+1}}{n^2}$, how would you bound the error after N terms?
  \item What role does monotonicity $(b_{n+1} \le b_n)$ play in the Alternating Series Test?
  \item Can an alternating series converge if the sequence $b_n$ is not decreasing? What does the test say about this?
  \item Compare the convergence behavior of $\sum \frac{(-1)^{n+1}}{n}$ and $\sum \frac{1}{n}$.
  \item Why is it often difficult to compute the exact sum of an alternating series?
  \item What practical advantage does the Alternating Series Estimation Theorem provide?
  \item How does the Alternating Series Test differ from comparison or ratio tests?
  \item When analyzing a series, why is it useful to check absolute convergence first?
  \item What happens if $\lim_{n \to \infty} a_n \neq 0$ for an alternating series?
\end{enumerate}
